\section{Structured Reader Replication}
\label{app:reader}
\begin{table*}[t]
\centering\small
\begin{tabular}{@{}lrrrrr@{}}
\toprule
 & \multicolumn{2}{c}{Primary} & \multicolumn{2}{c}{Diagnostic} & \\
Method & EM & F1 & EM & F1 & Valid\\
\midrule
Original BM25 & 23.33 & 24.44 & 9.09 & 9.09 & 71/71\\
Original + year & 25.00 & 25.00 & 18.18 & 18.18 & 71/71\\
Title-aware & 35.00 & 36.11 & 9.09 & 9.09 & 71/71\\
Earliest completion & 35.00 & 36.11 & 18.18 & 18.18 & 71/71\\
Midpoint completion & 35.00 & 36.11 & 9.09 & 9.09 & 71/71\\
Latest completion & 35.00 & 36.11 & 9.09 & 9.09 & 71/71\\
Interval control & 35.00 & 36.11 & 9.09 & 9.09 & 71/71\\
Possible support & 35.00 & 36.11 & 9.09 & 9.09 & 71/71\\
Guaranteed support & 35.00 & 36.11 & 18.18 & 18.18 & 71/71\\
Title-aware + year & 33.33 & 33.33 & 18.18 & 24.24 & 71/71\\
\bottomrule
\end{tabular}
\caption{Complete structured reader results. EM and strict answer-set F1 are percentages. Valid counts use question-method conditions. All conditions remain in the denominators.}
\label{tab:reader-structured}
\end{table*}


\paragraph{Unchanged evidence and scoring.}
The replication retains all 234 distinct requests from the earlier primary and diagnostic cohorts.
Every message, question, source excerpt, and method membership remains unchanged.
The ten methods produce 710 question--method conditions.
The model remains Qwen2.5-3B-Instruct with the same Q4\_K\_M weights and pinned llama.cpp runtime.
The strict parser, normalization, answer aliases, and article-cluster resampling also remain fixed.
Appendix~\ref{app:reader-legacy} retains the earlier decoding protocol and its complete results.

\paragraph{Structured decoding.}
A JSON schema constrains the output object and citation indices.
Answer strings remain unrestricted, and the schema does not prescribe the number of answers.
The completion budget increases to 512 tokens, with a 4,096-token context.
Temperature remains zero and the seed remains 1729.
Prompt caching stays disabled.
A tokenizer preflight confirms that every complete prompt fits this context.
The six lexicographically smallest request identifiers form a syntax pilot.
All six complete normally and pass the frozen parser before the full run proceeds.
No answer scores guide this gate.

The run initially uses one four-thread CPU server.
A recorded scheduling amendment adds a second independent server with identical settings.
Completed responses remain fixed; outstanding requests alternate between workers.
One interrupted request had no completed output and returns to the queue.
The resulting archive contains 234 new completed responses, including the six pilot responses.
The earlier 306 execution records remain unchanged.
Per-request hashes identify the exact inputs, decoder settings, weights, runtime, and raw response.

\paragraph{Context identity and answer cardinality.}
Possible and guaranteed support produce identical requests for every primary question.
Prompt deduplication therefore gives identical primary answers under both policies.
Their diagnostic requests differ for seven questions.
The primary targets contain one answer for 56 questions and two answers for four questions.
Every diagnostic target contains one answer.
The earlier primary decoder records no partial strict matches, which explains equality between exact match and strict F1.
The new audit reports target and prediction cardinalities separately for every method.

\paragraph{Complete outcomes.}
Every response stops normally and passes the frozen parser.
The longest completion contains 51 tokens.
Across unique requests, 75 predictions are empty, 157 contain one answer, and two contain two answers.
Empty predictions are valid abstentions and remain in every metric denominator.
Table~\ref{tab:reader-structured} reports all ten methods in both cohorts.

The title-aware primary row has 35.00\% exact match and 36.11\% strict F1.
Its partial match concerns Vladimir Bukovsky.
The prediction is \emph{Russian}; the annotated set contains \emph{British} and \emph{Russian}.
Strict set F1 therefore assigns $2/3$ credit while exact match assigns zero.
The same partial match occurs under each temporal policy.

The diagnostic possible--guaranteed comparison changes four answer sets, all with valid responses.
One change gains strict F1, and three retain zero strict F1.
Strict ties do not establish semantic equivalence.
For Mbaye-Jacques Diop, the prediction changes from the annotated party with its abbreviation to \emph{Socialist Party (PS)}.
Both receive zero under strict normalized matching.
The release records every changed question, answer, citation list, and cardinality.
The diagnostic retains its historical selection and nine article clusters.
Its guaranteed-minus-possible interval remains $[0.00,30.00]$ under the prespecified article bootstrap.
